\section{Naming and networking}
\label{sec:mobility}

Nothing in the pool configuration names a machine or an IP address. Concrete
values live outside version control.

\begin{table}[H]
\centering
\begin{tabularx}{\textwidth}{L{0.26\textwidth} L{0.26\textwidth} Y}
\toprule
\textbf{Layer} & \textbf{Value} & \textbf{Resolved how} \\
\midrule
Pool entry point & entry node's real hostname & mDNS, answered per interface \\
Principal namespace & \kn{<pool>.internal} & Never resolved; keeps IDTOKENS stable across naming changes \\
Host qualification & \kn{DEFAULT\_DOMAIN\_NAME = local} & Satisfies \kn{get\_full\_hostname()} \emph{and} is avahi-resolvable \\
Worker addresses & Advertised in each ClassAd & Condor dials them directly \\
\kn{FILESYSTEM\_DOMAIN} & Per-host default, never pool-wide & A shared value makes Condor skip file transfer entirely \\
\bottomrule
\end{tabularx}
\caption{Naming layers and how each is resolved.}
\end{table}

\subsection{Why names, not addresses}

Condor routes by the \textbf{advertised address}, the sinful string each daemon
puts in its ClassAd. It never resolves a name at connect time.

\begin{keypoint}
A worker that changes IP address is reachable again at its next collector update.
\textbf{DHCP churn on workers needs no DNS, no hosts file, and no action at all.}
Only the entry node must be findable by name, because every worker's
\kn{CONDOR\_HOST} points at it.
\end{keypoint}

\subsection{Why the pool must have a domain}

\kn{condor\_submit -remote} and \kn{condor\_q -name} call
\kn{get\_full\_hostname()}, which \textbf{rejects any name without a dot}.

A pool of bare hostnames cannot use remote submission at all. Since remote
submission is the entire point of this design, the domain is a hard requirement
rather than a tidiness preference.
\kn{DEFAULT\_DOMAIN\_NAME = local} qualifies every name and stays resolvable
through avahi via nsswitch (\kn{files mdns4\_minimal ... dns mdns4}).

\subsection{A multi-homed entry node serving two unroutable networks}

The pool spans two subnets that cannot route to each other, and machines move
between them. Both networks have ethernet and wifi, and a machine parks
wherever it happens to be. Plugging in a cable is not an option.

This works, using HTCondor's documented multi-homing support. An earlier revision
of this project's notes concluded it was impossible; that was wrong, and the
correction is the substance of this section.

\subsubsection{Why it looks impossible at first}

\kn{BIND\_ALL\_INTERFACES} is true and the daemons listen on \kn{0.0.0.0:9618},
so they \emph{accept} on every address. But a daemon advertises \textbf{one
primary address}, and peers are told to dial that one.

The result is a failure that mostly looks like success:

\begin{table}[H]
\centering
\begin{tabularx}{\textwidth}{L{0.30\textwidth} L{0.28\textwidth} Y}
\toprule
\textbf{Direction} & \textbf{Initiator} & \textbf{Result} \\
\midrule
worker $\rightarrow$ collector (register, update) & Worker dials \kn{CONDOR\_HOST} & Works, outbound \\
shadow $\rightarrow$ startd (claim) & Entry node dials the worker & Works \\
\textbf{starter $\rightarrow$ shadow (file transfer)} & \textbf{Worker dials the shadow's advertised address} & \textbf{Fails} \\
\bottomrule
\end{tabularx}
\caption{Three of four directions work, which is what makes this hard to see.}
\end{table}

\begin{trap}
The machine registers, jobs match, and jobs \textbf{run to completion}. Then the
return channel fails, the work is discarded, the job is re-queued, and it loops,
burning slots while \kn{condor\_status} shows a perfectly healthy machine.
\end{trap}

Two obvious fixes do not work, both tested:

\begin{itemize}
  \item \kn{NETWORK\_INTERFACE} with two addresses is accepted as a value, but
        \kn{addrs=} still carries one IPv4. The manual is explicit: \emph{``HTCondor
        daemons can only advertise two IP addresses\ldots One is the public IP
        address and the other is the private IP address.''}
  \item \kn{TCP\_FORWARDING\_HOST} set to a hostname is resolved \textbf{locally
        at startup} and baked in as a single IP. It also silently flipped the
        primary to the wifi address, which would have broken the wired machines
        instead.
\end{itemize}

Sinful strings are address-based by design: resolution happens once, at the
advertiser, and every peer receives the same string.

\subsubsection{What does work}

\kn{PRIVATE\_NETWORK\_NAME} and \kn{PRIVATE\_NETWORK\_INTERFACE}. The entry node
publishes \textbf{both} addresses; a peer uses the private one if it declares a
matching \kn{PRIVATE\_NETWORK\_NAME}, and the primary otherwise.

Two addresses is exactly enough for two networks.

\begin{lstlisting}
entry node:  PRIVATE_NETWORK_NAME      = $(FULL_HOSTNAME)
             PRIVATE_NETWORK_INTERFACE = <interface on the private segment>

advertises:  <PRIMARY:9618?PrivAddr=<PRIVATE:9618>
                          &PrivNet=<entry node hostname>&...>
\end{lstlisting}

The catch is that the choice is made by the \textbf{reader}. A hardcoded value on
a worker is silently wrong the moment that machine moves. In a fleet where
machines move constantly, that means it is wrong most of the time.

\subsubsection{Deriving it instead of configuring it}

\kn{PRIVATE\_NETWORK\_NAME} defaults to the machine's own \kn{FULL\_HOSTNAME},
which can never match another machine's. The default therefore already means
\emph{``use the primary address''}. Only machines on the private segment need to
do anything, and they can work it out for themselves:

\begin{lstlisting}[style=term]
1. resolve CONDOR_HOST      -> avahi answers per-interface, giving the
                               address reachable from HERE
2. ask the collector        -> its PrivAddr and its PrivNet name
3. reached it on PrivAddr?  -> adopt that PrivNet
   otherwise                -> write nothing; the default is already correct
\end{lstlisting}

\begin{keypoint}
The worker holds \textbf{no network fact, no address and no name}. It learns
everything from the entry node at startup. The same file goes on every machine,
and two workers on opposite segments reach opposite conclusions from identical
inputs.
\end{keypoint}

Wire it as an \kn{ExecStartPre} on \kn{condor.service} plus a NetworkManager
dispatcher hook that restarts Condor \textbf{only when the derivation changes},
so a routine lease renewal does not evict running jobs.

\subsection{The entry node needs the same hook for the opposite reason}

Workers derive. The entry node \emph{publishes}: it holds two addresses and lets
peers choose, so it has no derivation to watch. When one of its own addresses
changes, its daemons continue advertising the old one and nothing in the pool can
reach it.

\begin{trap}
Observed: the entry node's wifi roamed to another subnet, the collector went on
advertising an address that no longer existed, and nothing recovered until that
address happened to return. A worker-style hook, watching only the derivation,
would have been silent throughout.
\end{trap}

The hook therefore restarts on either trigger, a changed derivation or a changed
set of local addresses. A restart rather than a reconfig, because a daemon bakes
its advertised sinful string at startup and \kn{condor\_reconfig} does not
revisit it.

\textbf{Workers keep only the derivation trigger, deliberately.} A worker whose
address changes within its own segment needs no action: it is reachable again at
its next collector update. Adding the address trigger there would evict running
jobs to repair something that repairs itself. The entry node differs because
every worker dials \emph{it} by name, so a single stale address breaks every
inbound path at once.

\textbf{Fail safe:} any error, whether the name will not resolve, the collector
is unreachable or the ad is unparseable, must write \textbf{nothing} rather than
a stale name. Stale is precisely the failure this mechanism exists to prevent.

The derivation needs root, since querying the collector uses the daemon token in
\file{/etc/condor/tokens.d/}; running it as \kn{ExecStartPre} satisfies that.

\subsection{Why a role alias failed, and was retired}

The original design pointed configuration at the \emph{role} rather than at
whichever machine currently held it. Moving the entry node would mean starting an
alias service on the new box and stopping it on the old one, with no worker
changes at all. Two pieces were needed together:

\begin{itemize}
  \item \textbf{An mDNS alias}, fixing address lookups. avahi has no CNAME
        support, so a small service publishes \kn{avahi-publish -a} and
        republishes if the address changes. Publish with \kn{-R}: the real
        hostname already owns the reverse record, and a second one makes reverse
        lookups ambiguous, and that breaks Condor's security layer, not DNS.
  \item \textbf{\kn{NETWORK\_HOSTNAME}}, fixing name lookups. Without it the
        schedd still advertises its hardware name and \kn{-remote <alias>} fails
        with ``Can't find address of schedd''. \kn{SCHEDD\_NAME} cannot
        substitute: with no \kn{@} it becomes
        \kn{\$(SCHEDD\_NAME)@\$(FULL\_HOSTNAME)}.
\end{itemize}

\begin{trap}
\textbf{The alias was what broke multi-homing.} \kn{avahi-publish -a} binds a
name to \emph{one literal address} with no interface awareness, so the alias
answered every querier with the same address regardless of which network they
were on. Pointing \kn{NETWORK\_HOSTNAME} at it then forced the daemons to bind
and advertise that single address.
\end{trap}

A machine's \textbf{own} hostname behaves differently: avahi answers
per-interface, giving each querier the address it can actually reach. Verified
from both sides: the same real hostname resolved to the wired address from the
wired LAN and to the wifi address from wifi.

So \kn{CONDOR\_HOST} names the machine. The cost is losing role indirection:
relocating the entry node means updating \kn{CONDOR\_HOST} on every worker.
Correctness across networks beats the abstraction.

\begin{trap}
\textbf{Retiring a name means finding every consumer, not every worker.} When
the alias service was removed, the entry node and both workers were checked and
all three already named the real host. The submit-only client was not checked,
and it still named the alias, so remote submission from that machine broke and
stayed broken until it was next touched.
\end{trap}

That one failed loudly when finally exercised (\kn{Error: unknown host}), which
is the good case. A \emph{worker} left in that state would have registered
against nothing and looked healthy, which is the shape of every defect in
Chapter~\ref{sec:silent}. Grep every host for the name being retired, including
hosts that run no daemons.

\subsection{A VPN interface on a pool member is a hazard}
\label{sec:vpn}

Condor ranks candidate addresses by publicness and advertises the most public as
primary. Its private set is RFC1918, meaning \kn{10/8}, \kn{172.16/12} and
\kn{192.168/16}, so \textbf{tailscale's \kn{100.64/10} does not count as
private} and is a candidate to become the advertised address.

Measured on the first user PC, which has three interfaces:

\begin{table}[H]
\centering
\begin{tabularx}{\textwidth}{L{0.22\textwidth} L{0.26\textwidth} Y}
\toprule
\textbf{Interface} & \textbf{Address} & \textbf{In the pool?} \\
\midrule
wired & \kn{172.16.0.x/23} & yes \\
wifi & \kn{192.168.0.x/24} & yes \\
\kn{tailscale0} & \kn{100.x.x.x/32} & \textbf{no} \\
\bottomrule
\end{tabularx}
\caption{A pool member carrying an interface the pool cannot use.}
\end{table}

The entry node has no tailscale interface and cannot reach that address at all,
verified by connecting to each of the machine's three addresses from it.
Advertising the overlay address would produce the return-channel failure above:
the worker registers, jobs match, jobs run to completion, and then output
transfer fails and the work is discarded while \kn{condor\_status} shows a
healthy machine.

So \kn{NETWORK\_INTERFACE} is named explicitly on any machine carrying a VPN
interface, choosing the segment the entry node's \emph{primary} address is on so
that entry-to-worker takes the direct path. The cost is that the line is wrong
if the machine's attachment ever changes. That is accepted, because the
alternative is trusting a ranking that demonstrably picks wrong.

\subsubsection{The pin only constrains IPv4}

Verified afterwards on the slot ad, which is the only place this is visible:

\begin{lstlisting}[style=term]
<WIRED:9618?addrs=WIRED-9618
                 +[TAILSCALE-ULA-IPV6]-9618&...>
\end{lstlisting}

The primary is correct, but the \kn{addrs=} list still carries tailscale's
global IPv6. The pool interfaces have only link-local \kn{fe80::} addresses,
which Condor will not advertise, so the overlay's ULA was the only global IPv6
available, and it was selected independently of \kn{NETWORK\_INTERFACE}.

\begin{keypoint}
Nothing has broken: the entry node advertises IPv4 only and dials IPv4. But a
peer preferring IPv6 from that list would dial an address the entry node cannot
reach. \kn{ENABLE\_IPV6 = FALSE} closes it. \textbf{Check \kn{addrs=}, not
just the primary.}
\end{keypoint}

\subsection{Why UID\_DOMAIN is deliberately different}

\kn{UID\_DOMAIN} is a principal namespace and is never resolved as a hostname.
Keeping it distinct from \kn{DEFAULT\_DOMAIN\_NAME} means the naming scheme can
change without invalidating every token. Tokens carry
\kn{<user>@<pool>.internal}.

\begin{trap}
\textbf{Changing \kn{UID\_DOMAIN} invalidates every token.} Old tokens still
authenticate, but they map to a different principal, silently splitting
fair-share accounting rather than failing visibly.
\end{trap}

\subsection{Operational notes, learned the hard way}

\begin{itemize}
  \item Changing \kn{DEFAULT\_DOMAIN\_NAME} or \kn{NETWORK\_HOSTNAME} changes
        daemon names, so it needs a \textbf{restart}, not \kn{condor\_reconfig}.
  \item \textbf{Install configuration and restart in one step.} A daemon that
        fails to resolve the collector at startup does not retry until its next
        update interval, leaving ads missing for minutes.
  \item \kn{CONDOR\_HOST = \$(FULL\_HOSTNAME)} stops resolving the moment the
        domain changes, until that name resolves \emph{on the entry node itself}.
        Daemons come up but cannot reach their own collector.
  \item Hold code 16 (\kn{Spooling input data files}) is transient and normal
        during \kn{-spool} submission, not a failure.
  \item avahi can serve cached names for dead hosts, and \kn{getent} may return
        an IPv6 link-local address for a \kn{.local} name where avahi's IPv4
        answer is correct. Neither has broken anything yet; both are plausible
        sources of intermittent failures.
\end{itemize}

Vanilla-universe jobs need no worker-to-worker connectivity, because every
conversation is entry node to worker. The entry node is a star hub, not a router.
